\begin{textsample}{POS Dim 1 – df – Score 32.00 – df/df\_ef\_1.txt}  \label{ex:f1_pos_098}
APRESENTAÇÃO A Secretaria de Estado de Educação do Distrito Federal - SEEDF, reafirmando seu \textbf{compromisso} com \textbf{uma} educação de qualidade social para o sistema de ensino do Distrito Federal, e com o intuito de garantir \textbf{que} o currículo continue à serviço da aprendizagem de todos os estudantes, apresenta a 2ª edição do Currículo em Movimento do Distrito Federal para o Ensino Fundamental.
A sua 1ª edição, construção coletiva resultante de estudos e debates entre profissionais da educação, em seus pressupostos \textbf{teóricos}, assegura a identidade dinâmica do documento quando, ao se propor em movimento, prevê a necessidade de “ [ … ] ser permanentemente avaliado e significado a partir de concepções e práticas empreendidas por cada \textbf{um} e cada \textbf{uma} no contexto concreto das escolas e das salas de aula desta rede pública de ensino ” ( DISTRITO FEDERAL, 2014 ).
Após quatro anos de sua implementação, mesmo traduzido como \textbf{uma} referência para as redes de ensino no Distrito Federal, cujos \textbf{alicerces} epistemológicos corroboram \textbf{uma} educação \textbf{baseada} em \textbf{teorias} crítica e pós-crítica de currículo, a 1ª edição do Currículo em Movimento da Educação Básica \textbf{necessitava} de atualizações especialmente após a universalização da organização escolar em Ciclos para as Aprendizagens na rede pública de ensino em 2018.
Outra questão importante considerada para a revisitação desse documento foi \textbf{que}, com a homologação da Base Nacional Comum Curricular – BNCC em dezembro de 2017 ( Resolução CNE/CP nº 2 ), seguida de adesão da SEEDF ao Programa de Apoio à Implementação da BNCC, previsto na Portaria nº 331, do Ministério da Educação, surgiu a necessidade de alteração das matrizes curriculares a fim de contemplar os conhecimentos essenciais trazidos na BNCC, garantindo aos estudantes do Distrito Federal os mesmos direitos de aprendizagem assegurados a todos os outros estudantes brasileiros ( BRASIL, 2017 ).
No processo de construção da 2ª edição do Currículo para o Ensino Fundamental, a partir de discussões realizadas por professores de todos os componentes curriculares, como também das modalidades da Educação Básica, e diversos outros profissionais da educação, optou -se por manter as concepções \textbf{teóricas} e os princípios pedagógicos da 1ª edição do Currículo em Movimento : formação para Educação Integral ; Avaliação Formativa ; \textbf{Pedagogia} Histórico-Crítica e Psicologia Histórico-Cultural ; Currículo Integrado ; Eixos Integradores ( para os Anos Iniciais : Alfabetização, Letramentos e Ludicidade ; e, para os Anos Finais : Ludicidade e Letramentos ) e Eixos Transversais ( Educação para a Diversidade, Cidadania e Educação em e para os Direitos Humanos e Educação para a Sustentabilidade ).
Também primou -se pela manutenção da estrutura de objetivo de aprendizagem e conteúdo por entender \textbf{que} esses são elementos \textbf{que} corroboram os pressupostos \textbf{teóricos} assumidos \textbf{enquanto} \textbf{fundamentos} de currículo da SEEDF.
Importante \textbf{salientar} \textbf{que}, neste processo de visitação do Currículo para sua reelaboração, foram oportunizados, aos profissionais de educação e à sociedade civil, espaços para estudos, reflexões e discussões da proposta e contribuições diversas para elaboração final do documento : Fóruns Regionais ; Ciclo de Formações ; Ciclo de Plenárias por Componente Curricular e por Área de Conhecimento ; Leitores Críticos e Consulta Pública.
Houve ainda a constituição de \textbf{uma} Comissão Estadual, formalizada pela Portaria nº 163 da Secretaria de Estado de Educação do Distrito Federal, de 07 de junho de 2018, composta por representantes de dezesseis instituições : unidades administrativas da SEEDF, Poder Legislativo e entidades representativas de docentes e \textbf{discentes}, pública e privadas.
Esse grupo, de caráter consultivo, reuniu -se periodicamente com o papel de fortalecer o movimento, sugerindo melhorias ao processo.
Essencial também o destaque de \textbf{que} o material produzido foi de \textbf{um} currículo para a rede pública de ensino da SEEDF \textbf{que} sirva ainda de referencial curricular para a rede privada do Distrito Federal.
Esta, em respeito à sua forma de organização escolar – seja em ciclos ou em séries – fará os ajustes necessários ao documento, visando o desenvolvimento do seu trabalho.
Ao atender intrínseca necessidade de atualização do Currículo, a partir de contribuições de professores das redes de ensino e diversas entidades da sociedade civil, a SEEDF propõe como principais mudanças para a 2ª edição do Currículo em Movimento do Distrito Federal – Ensino Fundamental : - \textbf{um} único volume \textbf{que} contempla todo o Ensino Fundamental ( Anos Iniciais e Anos Finais ) a fim de \textbf{que} se tenha \textbf{uma} visão ampla do processo de aprendizagem dentro dessa etapa da Educação Básica ; - objetivos e conteúdos dispostos por ano, para viabilizar o trabalho das unidades escolares organizadas em série ; mas com traçado pontilhado \textbf{que} os separe dentro do mesmo bloco indicando possibilidade de transitarem no tempo proposto para esse bloco, o \textbf{que} oportuniza o trabalho das unidades escolares organizadas em ciclos ; - inserção dos conhecimentos essenciais trazidos pela BNCC não contemplados na edição anterior do Currículo em Movimento e/ou transferência dos objetivos e conteúdos para os anos em \textbf{que} eles são apresentados na Base ; - \textbf{contextualização} do Distrito Federal, ao ampliar elementos locais nos objetivos de aprendizagem ; - \textbf{abordagem} da cultura digital, explorando recursos midiáticos e características próprias de comunicação e informação ; - progressão dos objetivos de aprendizagem nos anos/blocos subsequentes a fim de \textbf{que}, gradualmente, ampliem -se e aprofundem -se os conhecimentos, minimizando assim os impactos ocorridos nos processos de transição entre os anos e inter e intrablocos.
Os objetivos de aprendizagem do Ensino Fundamental apresentados nas normativas pedagógicas da SEEDF, pautadas nas Diretrizes Curriculares Nacionais da Educação Básica – DCN ( 2013 ), visam : 1. possibilitar as aprendizagens, a partir da \textbf{democratização} de saberes, em \textbf{uma} perspectiva de inclusão considerando os Eixos Transversais : Educação para a Diversidade, Cidadania e Educação em e para os Direitos Humanos, Educação para a Sustentabilidade ; 2. promover as aprendizagens mediadas pelo pleno domínio da leitura, da escrita e do cálculo e a formação de atitudes e valores, permitindo vivências de diversos letramentos ; 3. oportunizar a compreensão do ambiente natural e social, dos processos histórico-geográficos, da diversidade étnico-cultural, do sistema político, da economia, da tecnologia, das artes e da cultura, dos direitos humanos e de princípios em \textbf{que} se fundamenta a sociedade brasileira, latino-americana e mundial ; 4. fortalecer vínculos da escola com a família, no sentido de proporcionar diálogos éticos e a corresponsabilização de papéis distintos, com vistas à garantia de acesso, permanência e formação integral dos estudantes ; 5. compreender o estudante como \textbf{sujeito} \textbf{central} do processo de ensino, capaz de atitudes éticas, críticas e reflexivas, comprometido com suas aprendizagens, na perspectiva do protagonismo estudantil.
Para \textbf{que} os estudantes alcancem os objetivos de aprendizagem, é fundamental \textbf{que} este Currículo seja vivenciado e reconstruído no cotidiano escolar, sendo, para tanto, \textbf{imprescindível} a organização do trabalho pedagógico da escola.
A utilização de estratégias \textbf{didático-pedagógicas} deve ser desafiadora e provocativa, levando em conta a construção dos estudantes, suas hipóteses e estratégias na resolução de problemas apresentados.
Como aspectos fundamentais para essa construção, constituem -se o Conselho de Classe, preferencialmente participativo ; a análise das aprendizagens para reorganização da prática docente ; a formação continuada no \textbf{lócus} da escola ; a coordenação pedagógica, como espaço e tempo primordiais para o trabalho coletivo ; entre outros. \textbf{Um} ambiente educativo com recursos variados, materiais didáticos atrativos e diversificados e situações problematizadoras \textbf{que} contemplem todas as áreas do conhecimento disponibilizadas aos estudantes são elementos capazes de promover as aprendizagens por meio da ação investigativa e criadora.
Para a qualificação da implementação deste Currículo nas unidades escolares, é essencial a articulação das diferentes áreas do conhecimento, com vistas à compreensão crítica e reflexiva da realidade. \textbf{Um} diálogo entre os conhecimentos, proposta \textbf{que} quebra a \textbf{fragmentação} do currículo na perspectiva coleção ( BERNSTEIN, 1977 ), demonstra \textbf{compromisso} ético no cumprimento da função social da escola.
A opção por \textbf{um} trabalho pautado nos princípios de unicidade teoria-prática, \textbf{interdisciplinaridade}, \textbf{contextualização} e flexibilização fortalece propósitos educacionais relevantes para a formação dos estudantes.
Nesse contexto, abre -se espaço para experiências, saberes, práticas dos \textbf{sujeitos} comuns \textbf{que} protagonizam e compartilham conhecimentos e vivências construídos em espaços sociais diversos.
Também dentro dessa perspectiva, os estudantes do Ensino Fundamental assumem, em seu percurso formativo, a condição de \textbf{sujeitos} de direito e constroem, gradativamente, sua cidadania ( BRASIL, 2013 ).
O trabalho pedagógico desenvolvido nas unidades escolares, portanto, deve estar voltado para as necessidades de aprendizagem de todos os estudantes, respeitando seus tempos de desenvolvimento, com a garantia de \textbf{um} processo contínuo de formação integral.
O ensino, então, não fica \textbf{restrito} à transmissão de conteúdos e à prática de avaliações \textbf{que} valorizem apenas o caráter quantitativo ao final de cada bimestre ; diferente disso, \textbf{aprimora} -se constantemente os processos de ensinar, de aprender e de avaliar, tendo como princípio fundamental a garantia das aprendizagens para todos os estudantes.
Visando \textbf{um} processo ininterrupto de aprendizagem, a compreensão de educação trazida neste Currículo adota o princípio da progressão continuada, \textbf{que} é \textbf{basilar} no modo de organização escolar em ciclos e pressupõe avanço nas aprendizagens dos estudantes, diferentemente da chamada promoção automática, caracterizada pela aprovação dos estudantes nos anos escolares independente da conquista das aprendizagens.
Como contribuição para \textbf{uma} educação transformadora da sociedade pela promoção das aprendizagens de todos os estudantes, alicerçada à perspectiva de \textbf{uma} avaliação encorajadora, esta Secretaria de Estado de Educação apresenta a 2ª edição do Currículo em Movimento do Distrito Federal para o Ensino Fundamental, como material passível de constante avaliação e alterações tendo em vista a necessidade de acompanhar inovações, estudos e discussões pedagógicas \textbf{atuais} tanto quanto aspectos da \textbf{contemporaneidade} \textbf{que} precisem ser considerados.

% matched lemmas: abordagem, alicerce, aprimorar, atual, basear, basilar, central, compromisso, contemporaneidade, contextualização, democratização, didático-pedagógico, discente, enquanto, fragmentação, fundamento, imprescindível, interdisciplinaridade, lócus, necessitar, pedagogia, que, restringir, salientar, sujeito, teoria, teórico, um, uma
\end{textsample}
